\subsection{A bounded quantum probability of the same finite action}
\label{sec:source-scalar-quantum}

The canonical commutator of the preceding execution is state independent.
A separate construction propagates the original action's vacuum covariance
and a coherent preparation, giving an actual bounded-effect probability.
It uses all $64$ modes of the same $q=5$ operator, without transferring the
spatial continuum error of another population.

Set $\hbar=1$, $Q=M^{1/2}u$, $B=M^{1/2}\mathsf A M^{-1/2}$ and
$[Q_i,P_j]=i\delta_{ij}$. On $L^2(\mathbb R^{64})$, let
\[
 H=\tfrac12(P^TP+Q^TBQ),\qquad
 \Phi(g)=(M^{1/2}g)^TQ,\qquad
 E_g=\tfrac12\bigl(I+\sin\Phi(g)\bigr).
\]
The effect satisfies $0\le E_g\le I$ by spectral functional calculus.
Since $B\ge I$, $H$ has a unique normalized ground state $\Omega$.
The baseline is $\Omega$; the intervention is
$\exp(i\Phi(v))\Omega$, with initial canonical means
$Q=0$, $P=M^{1/2}v$ and the same initial covariance.
Both states are propagated by the exact symmetric gates
\[
 S=e^{-i\tau Q^TBQ/4}e^{-i\tau P^TP/2}e^{-i\tau Q^TBQ/4}.
\]
At operator level the integrator memory is
$Q_{-1}=(I-\tau^2B/2)Q_0-\tau P_0$; the recorded preparation
fixes its mean, not a sharp quantum coordinate.
The reference uses $e^{-itH}$ on the same preparations.
The reference baseline is stationary; the split baseline is generally
squeezed. No modified-Hamiltonian vacuum is substituted.

\begin{proposition}[Original-vacuum probability comparison]
\label{prop:source-scalar-quantum-probability}
Write $r_*=30049/12005<3$, $\tau^2=1/245$,
$G_0=\langle g,g\rangle_M$ and
$G_1=\langle g,\mathsf A g\rangle_M$.
Let $V_0$ denote the vacuum variance of $\Phi(g)$ and $V_j$ its variance
after $j$ split steps, with or without the coherent displacement. Then
\begin{align}
 0\le V_j-V_0&\le
 D:=\frac{\tau^2\sqrt{G_0G_1}}{8(1-r_*/4)},
 &0\le V_0&\le G_0/2.\label{eq:scalar-quantum-variance}\\
 |p_j-p(t_j)|&\le \tfrac12\sqrt{G_0}\,e_j+\tfrac14D,
 &t_j&=j\tau,\label{eq:scalar-quantum-probability}
\end{align}
where $p_j$ and $p(t_j)$ are the intervention probabilities of $E_g$,
and $e_j$ is any valid enclosure of
$\|u_j-\mathsf A^{-1/2}\sin(t_j\sqrt{\mathsf A})v\|_M$.
Both baseline probabilities are exactly $1/2$, so the same estimate
bounds the difference of intervention-minus-baseline responses.
\end{proposition}
\begin{proof}
For a mode of frequency $\omega=\sqrt\lambda$, set
$x=\tau^2\lambda/4$, $\theta=2\arcsin(\tau\omega/2)$ and
$\widetilde\omega=\omega\sqrt{1-x}$. The split Heisenberg matrix is
\[
 \begin{pmatrix}1-2x&\tau\\-\tau\omega^2(1-x)&1-2x\end{pmatrix},
 \qquad
 Q_j=\cos(j\theta)Q_0+
       \frac{\sin(j\theta)}{\widetilde\omega}P_0.
\]
The original vacuum has variances $1/(2\omega)$ and $\omega/2$, with
zero symmetrized cross covariance. Hence, exactly,
\[
 \operatorname{Var}(Q_j)-\operatorname{Var}(Q_0)
 =\frac{\tau^2\omega\sin^2(j\theta)}{8(1-\tau^2\lambda/4)}\ge0.
\]
Summing with the squared spectral components of $M^{1/2}g$ gives the
first bound in \eqref{eq:scalar-quantum-variance}, using
$\langle g,\sqrt{\mathsf A}g\rangle_M\le\sqrt{G_0G_1}$.
The inequality $\mathsf A\ge I$ gives $V_0\le G_0/2$.
For a Gaussian state with field mean $m$ and variance $V$,
\[
 \langle E_g\rangle=\tfrac12\bigl(1+e^{-V/2}\sin m\bigr).
\]
The displacement changes only the mean; its split mean is
$\langle g,u_j\rangle_M$.
Cauchy--Schwarz bounds the mean error by $\sqrt{G_0}e_j$.
The sine is $1$-Lipschitz and $e^{-V/2}$ is $1/2$-Lipschitz for
$V\ge0$, proving \eqref{eq:scalar-quantum-probability}.
Centered Gaussian states have zero expectation of the odd sine,
including the squeezed split baseline.
The mode calculations hold on the preserved Schwartz domain;
the final effect is bounded on the full Hilbert space.
\end{proof}

\paragraph{Exact finite probability enclosures.}
The mass and energy moments are exact elements of $\mathbb Q(\phi)$.
An independent rational verifier consumes the authenticated execution and
its full-mode mean-error bounds at all $21$ stored positive times.
It obtains the uniform outward bound $V_j-V_0\le0.000646532908$.
No covariance diagonalization is needed: rational bounds on the moments,
$e^{-z}\ge1-z$, and $\sin y\ge y-y^3/6$ for $0\le y\le1$
give signed probability intervals. All times are retained; at $15$ of them
the certified response exceeds the total comparison error.
For example, at $t=16\sqrt5/35$,
\[
 p_{16}\in[0.490342191369,\,0.490493790733],\qquad
 |p_{16}-p(t)|\le0.000810888870.
\]
Thus the split response to the baseline $1/2$ has magnitude at least
$0.009506209267$, and the continuous finite-action response has the same
negative sign with magnitude at least $0.008695320397$.

The action, boundary, canonical quantization, original vacuum, coherent
preparation, detector normalization and model tick are explicit inputs.
These are bounded quantum effect predictions at executed finite-action
times. They do not attest observed outcomes, a quantum state encoded by
the classical registers, a spatial continuum error, intermediate-time
accuracy, an interacting quantum field, native time or laboratory units.
